\section{Lattice fields}\label{sec:uv-lattice}

The lattice measure and the local observables are fixed together, before any sequence of volumes is prescribed.
This section proves the exact symmetries of the lattice law and defines the classical labels of local observables.
The normalization of these polynomials has two stages: their Taylor coefficients reproduce classical fields, and their coefficients are subsequently changed by the renormalization transformation.
Keeping these stages distinct will allow us to use one finite stencil at every cutoff.

Throughout this section, $G\subset U(m)$ is the closed matrix group of the theorem, with real Lie algebra~$\g$.
The arguments here do not require simplicity or connectedness.
The Frobenius norm on matrices is denoted by~$\|\cdot\|_{\mathrm F}$.

\subsection{The Wilson measure}

Only positive links are independent variables.
For $\ell\ge1$, put $\Lambda_\ell=(\Z/\ell\Z)^4$ and assign $U(y,\mu)\in G$ to each link from~$y$ to~$y+e_\mu$.
Transport on the reverse link is $U(y,\mu)^{-1}$.
If $w$ is a lattice walk from~$x$ to~$z$, write $U(w)$ for the ordered product of its link variables.
A gauge transformation $\gamma:\Lambda_\ell\to G$ acts by
\begin{equation}\label{eq:uv-gauge}
 U^\gamma(y,\mu)=\gamma(y)U(y,\mu)\gamma(y+e_\mu)^{-1},
 \qquad U^\gamma(w)=\gamma(x)U(w)\gamma(z)^{-1}\,.
\end{equation}
The second identity follows by cancellation at the intermediate vertices, including for backward steps.

For~$u\in G$, define
\begin{equation}\label{eq:uv-chord}
 s(u)=1-\frac1m\operatorname{Re}\tr u,
 \qquad \chi(u)=\frac{\|u-1\|_{\mathrm F}}{\sqrt m}\,.
\end{equation}
Unitarity and the trace identity for the Frobenius norm give
\begin{equation}\label{eq:uv-chord-square}
 \chi(u)^2=2s(u),\qquad 0\le s(u)\le2,\qquad
 \chi(uv)\le\chi(u)+\chi(v)\,.
\end{equation}
Indeed, $\|u-1\|_{\mathrm F}^2=2m-2\operatorname{Re}\tr u$, and $\|u-1\|_{\mathrm F}\le2\sqrt m$.
For the last inequality write $uv-1=(u-1)+u(v-1)$ and use invariance of the Frobenius norm under unitary multiplication.
Both~$s$ and~$\chi$ are continuous and invariant under conjugation and inversion.

Let~$dh$ be probability Haar measure on~$G$, and let~$d\eta_\ell$ be its product over the~$4\ell^4$ positive links.
Write
\begin{equation}\label{eq:uv-wilson}
 \begin{aligned}
 U_{y;\mu\nu}
 &=U(y,\mu)U(y+e_\mu,\nu)U(y+e_\nu,\mu)^{-1}U(y,\nu)^{-1},\\
 E_\ell(U)&=\sum_{y\in\Lambda_\ell}\sum_{\mu<\nu}s(U_{y;\mu\nu}),\\
 Z_\ell(\beta)&=\int e^{-\beta E_\ell(U)}\,d\eta_\ell(U)\,.
 \end{aligned}
\end{equation}
Each site and increasing axis pair contributes once, including when the periodic side is small.
There is no further factor in this definition of the energy.

Let~$H_4$ be the group of signed permutation matrices.
For~$A(y)=hy+z$, with~$h\in H_4$ and~$z\in\Lambda_\ell$, the pullback~$A^*U$ assigns to a positive link its image transport, using the inverse variable if the image direction is negative.

\begin{proposition}[The lattice law]\label{prop:uv-law}
For every~$\ell\ge1$ and~$\beta\in\R$, the measure
\begin{equation}\label{eq:uv-law}
 d\mu_{\ell,\beta}(U)=Z_\ell(\beta)^{-1}e^{-\beta E_\ell(U)}\,d\eta_\ell(U)
\end{equation}
is a Borel probability measure on the compact configuration space.
Every continuous complex function and every bounded Borel complex function is integrable.
Every gauge transformation and every affine signed permutation~$A$ preserve this measure.
Consequently, for every integrable complex~$F$,
\begin{equation}\label{eq:uv-expectation-symmetry}
 \int F(U^\gamma)\,d\mu_{\ell,\beta}(U)
 =\int F(U)\,d\mu_{\ell,\beta}(U)
 =\int F(A^*U)\,d\mu_{\ell,\beta}(U)\,.
\end{equation}
\end{proposition}

\begin{proof}
The group~$G$ is compact, as a closed subgroup of the compact unitary group.
Thus the finite product of link spaces is compact.
The energy is continuous and satisfies~$0\le E_\ell\le12\ell^4$, by~\eqref{eq:uv-chord-square} and the six increasing axis pairs.
The density is therefore bounded above and bounded below by a positive constant, for every real~$\beta$.
Its integral~$Z_\ell(\beta)$ is finite and positive, and division by that integral proves normalization.
Continuous functions on the compact configuration space are bounded; the integrability assertions follow.

Probability Haar measure on a compact group is invariant under right translations and inversion as well as left translations.
For example, the image of Haar measure under a fixed right translation is again left invariant with total mass one, so Haar uniqueness gives equality.
Inversion then exchanges left and right invariance and also preserves the normalization.
Each gauge transformation acts on a link by fixed left and right multiplications.
It preserves product Haar measure, while~\eqref{eq:uv-gauge} makes every plaquette transform by conjugation at its initial vertex.
Hence it preserves~$E_\ell$ and the density in~\eqref{eq:uv-law}.

An affine signed permutation permutes the positive-link indices and inverts selected link variables.
The shift of the initial vertex on a reversed link is included in this permutation.
It therefore preserves product Haar measure.
The image of an elementary square is another elementary square; changing its initial corner conjugates its transport, and reversing its circulation inverts it.
To check the indices, write~$he_\mu=\sigma_\mu e_{\pi(\mu)}$, with~$\sigma_\mu\in\{-1,1\}$.
The image square has its positive-axis base at
\begin{equation*}
 hy+z-\mathbf1_{\{\sigma_\mu=-1\}}e_{\pi(\mu)}
      -\mathbf1_{\{\sigma_\nu=-1\}}e_{\pi(\nu)}
\end{equation*}
and its axis pair is the increasing rearrangement of~$\{\pi(\mu),\pi(\nu)\}$.
The axis pairs are permuted, and for each pair the site map is an affine bijection.
Thus the face indices are permuted, including on a periodic lattice of side one or two.
Conjugation and inversion invariance of~$s$ give invariance of~$E_\ell$.
Changing variables in~\eqref{eq:uv-law} proves~\eqref{eq:uv-expectation-symmetry}.
\end{proof}

\subsection{Classical labels}

The local polynomial labels describe a connection in radial gauge near one point.
They are defined before imposing field equations or taking a continuum limit.
For~$j\ge0$, let~$\mathcal A_j$ be the real vector space of homogeneous polynomial one-forms of degree~$j+1$ with values in~$\g$, satisfying
\begin{equation}\label{eq:uv-radial}
 \sum_{\mu=1}^4 z^\mu A_{j,\mu}(z)=0\qquad(z\in\R^4)\,.
\end{equation}
Equivalently, each component is the diagonal evaluation of a symmetric~$(j+1)$-linear map, with the same constraint.
These are Taylor coefficients, so their definition includes the factorial normalization of a Taylor polynomial.
Let~$J=(A_j)_{j\ge0}$ denote a sequence of such coefficients.
Choose real coordinates on each~$\mathcal A_j$, give every coordinate weight~$j+2$, and let~$V_d$ be the complex polynomials of total weight at most~$d$ invariant under simultaneous conjugation by~$G$.

The ordinary polynomial degree will count factors, separately from weight.
An orthogonal matrix acts on one-forms by
\begin{equation}\label{eq:uv-label-action}
 (T_hA_j)_\mu(z)=\sum_\nu h_{\mu\nu}A_{j,\nu}(h^{-1}z),
 \qquad (h\cdot v)(J)=v(T_{h^{-1}}J)\,.
\end{equation}
The action preserves the radial condition and commutes with conjugation by~$G$.
Complex conjugation of polynomial coefficients is denoted by~$v\mapsto\bar v$.

\begin{lemma}[Finite label spaces]\label{lem:uv-labels}
For every~$d\ge0$, the space~$V_d$ is finite dimensional and its elements are determined by their values on real jets.
It contains~$1$ and is naturally included in~$V_{d+1}$.
The decomposition by weight and factor number preserves~$V_d$, and its components are preserved by~$O(4)$ and complex conjugation.
There exist bases in all degrees, nested as sets of polynomials, consisting of real polynomials homogeneous in both gradings and containing~$1$.
\end{lemma}

\begin{proof}
Only coordinates with~$j+2\le d$ can occur in a monomial of weight at most~$d$.
There are finitely many such coordinates, and the exponent of each is bounded by~$d$.
Thus finitely many monomials span all polynomials under consideration, and their invariant subspace is finite dimensional.
A complex polynomial vanishing on every real jet vanishes identically, by the one-variable polynomial identity theorem applied successively in these finitely many real coordinates.

Conjugation by~$G$ acts linearly within each~$\mathcal A_j$.
It therefore preserves both weight and factor number, so comparison of homogeneous components of the invariance identity shows that every component remains invariant.
The same observation applies to~\eqref{eq:uv-label-action}.
All these real linear changes commute with coefficient conjugation.
Each invariant homogeneous component is the complex span of its real polynomials, by decomposition into real and imaginary coefficient parts.
Choose a real basis in each component, starting with the basis~$\{1\}$ of the component with no factors, and retain previous choices as~$d$ increases.
\end{proof}

Fix such bases and write~$b_{d,i}$ for their elements, with weights~$w_{d,i}$ and factor numbers~$f_{d,i}$; take~$b_{d,0}=1$.
We identify a polynomial with its coefficient vector in this basis when applying linear maps to labels.
The finite-dimensional continuum spaces~$F_d$ in the theorem will be quotients of labels after a limiting correlation table has been obtained.
No such quotient is taken here.

\subsection{Local representatives}

A classical label is a polynomial on Taylor coefficients; the lattice insertion must instead be a function of finitely many links.
We specify the normalization relating the two before constructing such representatives in the next section.

For a finite set~$\Sigma\subset\Z^4$, a function is local on~$\Sigma$ if it depends only on positive links whose two endpoints lie in~$\Sigma$.
A polynomial on these links means a polynomial in their matrix entries and conjugate entries.
It is continuous, since the coordinate functions are continuous.
At~$o\in\Lambda_\ell$, place a cover function by substituting~$U(o+[x]_\ell,\mu)$ for every cover link~$(x,\mu)$.
This substitution is meaningful even if the stencil wraps around the torus or distinct cover links become the same torus link.

To state the classical normalization, truncate a real jet~$J$ to weights at most~$d$ and write~$A^J_\mu=\sum_{j+2\le d}A_{j,\mu}$.
For a spacing parameter~$\varepsilon\in\R$, define its link transport by
\begin{equation}\label{eq:uv-link-ode}
 \frac{d}{dt}u(t)=u(t)\varepsilon A_\mu^J(\varepsilon(x+te_\mu)),\qquad
 u(0)=1,\qquad U_{J,\varepsilon}(x,\mu)=u(1)\,.
\end{equation}
This convention fixes the order of multiplication in the link transport.
For completeness, successive integration gives a normally convergent series on compact parameter sets: the term with~$r$ integrations is bounded by~$M^r/r!$ if the coefficient matrix is bounded by~$M$ on~$[0,1]$.
The series solves the equation, and iteration of the integral equation proves uniqueness.
In particular, the endpoint is real analytic in~$\varepsilon$ near~$0$, for each fixed real jet.
Approximating the coefficient by step functions gives convergent products of exponentials from~$\g$; their limits lie in the closed group~$G$.
Thus these are actual group-valued transports.

\begin{definition}[Normalized local section]\label{def:uv-sections}
A normalized local section in degree~$d$ consists of a finite~$H_4$-invariant set~$\Sigma_d\subset\Z^4$ containing~$0$ and a complex linear map~$v\mapsto P_d[v]$ to polynomials on links with both endpoints in~$\Sigma_d$.
Every~$P_d[v]$ is gauge invariant, and
\begin{equation}\label{eq:uv-section-symmetry}
 P_d[1]=1,\qquad P_d[\bar v](U)=\overline{P_d[v](U)},\qquad
 P_d[v](h^*U)=P_d[h\cdot v](U)\quad(h\in H_4)\,.
\end{equation}
For every real jet~$J$ and every~$0\le k\le d$, the section satisfies
\begin{equation}\label{eq:uv-section-taylor}
 \left.\frac{d^k}{d\varepsilon^k}P_d[v](U_{J,\varepsilon})\right|_{\varepsilon=0}
 =\left.\frac{d^k}{d\varepsilon^k}v(\operatorname{sc}_\varepsilon J)\right|_{\varepsilon=0},
 \qquad (\operatorname{sc}_\varepsilon J)_j=\varepsilon^{j+2}A_j\,.
\end{equation}
\end{definition}

The Taylor identity is an equality of classical coefficients at zero spacing, not an equality of observables at nonzero spacing.
No compatibility of~$P_d$ and~$P_{d+1}$ is imposed here.
Existence of these sections, with the derivative bounds needed for blocking, is a separate construction.
The conditions above alone do not assert it.

\subsection{Renormalized insertions}

To define a renormalized insertion, one needs an invertible linear map on the entire finite-dimensional label space.
It may mix different weights up to~$d$, including the constant label.
We specify here how this map enters the lattice observable; its construction and estimates belong to the subsequent ultraviolet sections.

Let~$\theta_N>0$ be the finest-level small-field threshold, and let~$S_t$ multiply the basis polynomial~$b_{d,i}$ by~$t^{f_{d,i}}$.
For an invertible complex linear map~$\tau_{N,p,d}$ on~$V_d$, define
\begin{equation}\label{eq:uv-insertion}
 I_{N,p}((d,v),o;U)
 =a_N^{-d}P_d[S_{\theta_N^{-1}}\tau_{N,p,d}^{-1}v](o;U)\,.
\end{equation}
The map~$\tau_{N,p,d}$ reads the running coefficients at a reference level~$r(d)$ fixed independently of~$N,p$.
More precisely, if~$x_m(v)$ are the linear coefficient maps for an insertion transported from level~$N$, their endpoint conditions are
\begin{equation}\label{eq:uv-reference-coordinates}
 x_N(v)=v,\qquad x_{r(d)}(v)=\tau_{N,p,d}v,\qquad
 x_m(a_N^d1)=a_m^d1\quad(r(d)\le m\le N)\,.
\end{equation}
These conditions concern coefficients of the inserted density, rather than pointwise values of the coarse observable.

\begin{proposition}[Locality of the insertions]\label{prop:uv-insertions}
Given normalized local sections, positive thresholds and invertible maps as in~\eqref{eq:uv-insertion}, the insertions are complex linear in~$v$, continuous, gauge invariant and local on~$o+[\Sigma_d]_\ell$, for~$\ell=L_pD_N$.
They satisfy the literal translation identity
\begin{equation}\label{eq:uv-insertion-translation}
 I_{N,p}((d,v),o;U(\cdot+z))=I_{N,p}((d,v),o+z;U)\,.
\end{equation}
If~$N\ge r(d)$, the maps~$\tau_{N,p,d}$ commute with conjugation and~$H_4$, and~\eqref{eq:uv-reference-coordinates} holds, then the insertions respect reality and, for~$A(o)=ho+z$,
\begin{equation}\label{eq:uv-insertion-covariance}
 I_{N,p}((d,h\cdot v),A(o);U)=I_{N,p}((d,v),o;A^*U),
 \qquad I_{N,p}((d,a_{r(d)}^d1),o;U)=1\,.
\end{equation}
Their running coefficients at level~$r(d)$, with terminal input~$\tau_{N,p,d}^{-1}v$, equal~$v$.
\end{proposition}

\begin{proof}
Every assertion in the first sentence follows by composition of the linear coefficient maps with the polynomial section of Definition~\ref{def:uv-sections}.
Periodic placement reads only links with both endpoints in the image of~$\Sigma_d$; identifying distinct cover vertices does not enlarge this set.
Translating the placement proves~\eqref{eq:uv-insertion-translation}.
Factor scaling commutes with conjugation and with~$H_4$, since these actions preserve factor number.
Their commutation with~$\tau_{N,p,d}$ and~\eqref{eq:uv-section-symmetry} give reality and affine covariance.
The reference endpoint and unit condition yield~$\tau_{N,p,d}(a_N^d1)=a_{r(d)}^d1$.
Since~$S_t1=1$ and~$P_d[1]=1$, substitution in~\eqref{eq:uv-insertion} proves its unit identity.
Finally, $x_{r(d)}(\tau_{N,p,d}^{-1}v)=v$ by~\eqref{eq:uv-reference-coordinates}.
\end{proof}

For the token set~$\mathcal T=\bigsqcup_dV_d$, Proposition~\ref{prop:uv-insertions} supplies, under its hypotheses, the continuity, gauge invariance and finite-stencil assertions Reg2--Reg4.
They hold at every cutoff; the additional reference normalization is asserted only once~$N\ge r(d)$.
With integrability supplied by Proposition~\ref{prop:uv-law}, the lattice moments and smeared moments in Part~I are therefore ordinary normalized integrals of these polynomial observables.

\gap{The existence of normalized local sections with their required derivative bounds, and the same-family construction of the invertible transfer maps satisfying~\eqref{eq:uv-reference-coordinates}, remain to be proved in the later ultraviolet sections. The local insertion proposition is conditional on these inputs.\sourcefootnote{Local sections: \texttt{YM/Integrated.lean}, lines 214--219. Transfer conditions: \texttt{YM/Source/Insertion.lean}, lines 89--114; same-family application: \texttt{YM/FullVolume/Cross.lean}, lines 30--71}}
